\section{ADFTD Split-3 Methodology (Supplementary)}
\label{app:adftd_split}

The ADFTD cohort is natively single-recording-per-subject, and our
nested variance decomposition (Sec.~\ref{sec:methods_variance})
requires multiple observations per subject to identify the
within-subject variance component. This appendix documents the
split-3 protocol, its role in the decomposition, and one semantic
caveat that should be held in mind when comparing variance fractions
across datasets.

\paragraph{Protocol.} Each subject's original eyes-closed resting
recording ($\sim\!14$ minutes, 500\,Hz, 19 channels at standard
10--20 positions) is partitioned into $3$ non-overlapping contiguous
segments of $\sim\!4.5$ minutes each. The segmentation is strictly
temporal: segment $1$ is the first third of the raw recording, segment
$2$ the second, segment $3$ the third. No epoch-level shuffling is
performed before segmentation, so segment boundaries do not mix data
from different time periods. All three segments of a subject carry
the same AD/HC diagnostic label. Subject-level cross-validation
groups all three segments of a subject into the same fold, so no
label leakage is introduced.

\paragraph{Why the split is methodologically necessary.} The pooled
variance decomposition used throughout Sec.~\ref{sec:results_atlas}
partitions frozen-FM representation variance into a subject component
$\sigma^2_\text{subj}$, a label-within-subject component
$\sigma^2_\text{label|subj}$, and a residual
$\sigma^2_\text{res}$ via a two-level nested ANOVA on $(n_\text{obs},
d_\text{feat})$ feature matrices. With one observation per subject,
$\sigma^2_\text{subj}$ and $\sigma^2_\text{res}$ are mutually
unidentifiable: each subject is represented by a single point in
feature space, and no estimate of the replicate spread around that
point is available. Multiple observations per subject are therefore
a prerequisite, not a methodological choice. The other three
datasets used in the variance-decomposition analysis (UCSD Stress,
EEGMAT, TDBRAIN) natively contribute multiple recordings per subject
and require no analogous splitting.

\paragraph{The split is not used for classification.} ADFTD
classification numbers reported in Sec.~\ref{sec:results_fm_adds}
use global prediction pooling: all three segments of a subject
contribute to a single per-subject prediction, and balanced accuracy
is computed at the subject level. The split-3 structure never
inflates the effective classification $N$ --- it is strictly a
device for making the within-subject variance component in
Sec.~\ref{sec:results_atlas}'s ANOVA identifiable.

\paragraph{Semantic-asymmetry caveat.} ADFTD within-subject variance
estimated from temporal segments of a single recording is not
semantically identical to within-subject variance on UCSD Stress or
EEGMAT, where subjects contribute multiple separate sessions (or a
rest-vs-task pair). The ADFTD estimate captures short-timescale
temporal fluctuation within a single seated 14-minute recording;
the Stress/EEGMAT estimates capture across-session or across-task
variability that includes setup, alertness, and task-structure
differences. Any literal numeric comparison of
$\sigma^2_\text{within}$ fractions across datasets with different
observation structures should therefore be read with this semantic
asymmetry in mind. In our reporting we use the ratio
$\sigma^2_\text{label|subj} / \sigma^2_\text{subj}$ as the primary
SDL-relevant metric, which is within-dataset and does not require
cross-dataset comparability of the residual component. The
qualitative subject-dominance ordering (subject variance
$\gg$ label variance) holds across all four datasets;
quantitative magnitudes of variance fractions should not be read as
directly commensurable across datasets.

\paragraph{Direction of any bias.} If the autocorrelation between
ADFTD temporal segments makes within-subject variance smaller than
it would be across separate sessions, the effect is a \emph{smaller}
denominator in the subject-dominance ratio, and therefore a
\emph{larger} ratio. Any bias introduced by the split-3 structure
thus works against the conservative reading of our ADFTD
subject-dominance claim, not in favour of it.
